\section{\MakeUppercase{Методика дослідження}}
\label{sec:chapter1}

\subsection{Стратегія пошуку та формування вибірки}
\label{sec:search_strategy}

Бібліометричний аналіз є визнаним методом кількісного дослідження наукової літератури~\cite{Pritchard1969, Donthu2021}, що дозволяє виявити структуру дослідницького поля, ідентифікувати ключові теми та простежити їх динаміку~\cite{Zupic2015}. Серед основних підходів бібліометричного аналізу --- ко-цитатний аналіз~\cite{Small1973}, аналіз ко-входжень ключових слів~\cite{Callon1983} та бібліометричне картографування~\cite{Cobo2011}. Сучасний огляд методів та інструментів бібліометричного аналізу наведено в~\cite{Aria2017}.

Джерелом бібліометричних даних обрано базу Web of Science Core Collection (WoS), яка є однією з найавторитетніших мультидисциплінарних наукометричних баз даних~\cite{Birkle2020}. Пошук публікацій здійснювався за запитом, спрямованим на виявлення досліджень моніторингу у сфері професійно-технічної освіти:

\begin{verbatim}
TS=("vocational education" OR "vocational training" OR
    "VET" OR "TVET") AND TS=("monitoring" OR
    "assessment" OR "evaluation" OR "quality assurance")
\end{verbatim}

Результати пошуку обмежувались англомовними публікаціями за період 1997--2025 років. Було отримано 1998 записів, які експортовано у форматі WoS (повні записи з цитованими посиланнями). Через технічне обмеження WoS на експорт (максимум 500--1000 записів за один раз), дані завантажувались порціями та об'єднувались у файл \texttt{merged\_data.txt} (1998 записів із повними полями, включаючи DE --- авторські ключові слова та ID --- Keywords Plus). Окремо файл \texttt{savedrecs.txt} містить 1000 записів без полів DE/ID, який використовувався для аналізу в~VOSViewer.

Таким чином, вибірка включає 1998 публікацій для описової статистики (аналіз Python) та 1000 публікацій для текстового аналізу та побудови бібліометричних карт (аналіз VOSViewer).

\subsection{Аналіз засобами VOSViewer}
\label{sec:vosviewer}

VOSViewer --- це програмний інструмент для побудови та візуалізації бібліометричних мереж, розроблений Н.\,Я.~ван Еком та Л.~Вальтманом у Лейденському університеті~\cite{vanEck2010}. Програма дозволяє створювати мережі на основі ко-цитування, бібліографічного зв'язку та ко-входжень ключових слів~\cite{vanEck2014}. VOSViewer використовує метод VOS (Visualization of Similarities) для розташування елементів на двовимірній карті та алгоритм кластеризації на основі модулярності для групування пов'язаних елементів~\cite{Waltman2010}. Для кластеризації за ко-цитуванням програма також підтримує інтеграцію з CitNetExplorer~\cite{vanEck2017}.

Для побудови карти ключових слів у VOSViewer було виконано текстовий аналіз (text mining) файлу \texttt{savedrecs.txt} із такими параметрами:
\begin{itemize}
    \item метод підрахунку: бінарний (binary counting) --- враховується лише факт наявності терміна в документі;
    \item мінімальна кількість входжень терміна: 100;
    \item після автоматичного відбору отримано 31 термін, з яких 3 нерелевантні (\textit{article}, \textit{patient}, \textit{year}) було виключено вручну;
    \item застосовано тезаурус для об'єднання синонімічних та морфологічних варіантів: усі терміни, що стосуються ПТО, позначено як <<vet>>, а терміни, що стосуються тих, хто навчається --- як <<student>>.
\end{itemize}

У результаті отримано карту з 28 ключових слів, розподілених на 4 кластери:

\begin{enumerate}
    \item \textbf{Кластер~1 <<Системи та політика ПТО>>} (11 термінів: vet, training, development, model, paper, country, company, field, order, germany, higher education) --- макрорівень, інституційний та системний вимір ПТО.
    \item \textbf{Кластер~2 <<Емпіричні дослідження результатів навчання>>} (8 термінів: student, study, school, assessment, performance, effect, group, difference) --- мікрорівень, кількісні дослідження учнів та вимірювання результатів.
    \item \textbf{Кластер~3 <<Педагогічна практика та компетентнісний підхід>>} (6 термінів: teacher, competence, practice, learning, work, context) --- мезорівень, взаємодія викладача з учнями та організація навчального процесу.
    \item \textbf{Кластер~4 <<Когнітивні результати та навички>>} (3 терміни: skill, knowledge, problem) --- мікрорівень, цільові результати ПТО.
\end{enumerate}

Просторове розташування кластерів на карті VOSViewer відображає семантичну структуру дослідницького поля: горизонтальна вісь інтерпретується як <<Система~$\leftrightarrow$~Особистість>>, а вертикальна --- як <<Результати/Вимірювання~$\leftrightarrow$~Процес/Контекст>>.

\subsection{Аналітична система Python}
\label{sec:python_system}

Для порівняння з результатами VOSViewer розроблено аналітичну систему на мові Python~3, що складається з 9 скриптів:

\begin{enumerate}
    \item \textbf{\texttt{wos\_parser.py}} --- спільний модуль парсингу файлів WoS. Реалізує функції \texttt{parse\_wos()}, \texttt{load\_savedrecs()}, \texttt{load\_merged()}, \texttt{load\_vosviewer\_map()} для читання даних у різних форматах.
    \item \textbf{\texttt{descriptive\_stats.py}} --- описова статистика вибірки (1998 записів): розподіл публікацій за роками, типами документів, джерелами, авторами, країнами, авторськими ключовими словами та Keywords Plus. Генерує 7 діаграм.
    \item \textbf{\texttt{text\_mining.py}} --- відтворення текстового аналізу VOSViewer: токенізація абстрактів та заголовків, бінарний підрахунок, фільтрація за порогом $\geq$100, застосування тезаурусу, збереження порівняльної таблиці частот.
    \item \textbf{\texttt{cooccurrence.py}} --- побудова матриці ко-входжень 28 ключових слів (28$\times$28), обчислення Total Link Strength (TLS), порівняння з~VOSViewer.
    \item \textbf{\texttt{clustering.py}} --- кластеризація мережі алгоритмом Лувена~\cite{Blondel2008} із перебором параметра resolution. Оцінка якості за метриками ARI~\cite{Hubert1985} та NMI порівняно з кластерами VOSViewer.
    \item \textbf{\texttt{network\_viz.py}} --- візуалізація мережі ключових слів: за координатами VOSViewer, за spring layout NetworkX~\cite{Hagberg2008}, та порівняння <<пліч-о-пліч>>.
    \item \textbf{\texttt{overlay\_viz.py}} --- накладні карти (overlay): середній рік публікації, середня цитованість, середня нормалізована цитованість.
    \item \textbf{\texttt{density\_viz.py}} --- карта щільності ключових слів на основі ядерної оцінки щільності (KDE).
    \item \textbf{\texttt{comparison.py}} --- підсумкова таблиця порівняння метрик VOSViewer та Python (CSV та LaTeX), обчислення кореляцій Пірсона, текстовий звіт українською мовою.
\end{enumerate}

Використані бібліотеки Python: pandas~\cite{McKinney2010} та numpy --- для обробки даних; NetworkX~\cite{Hagberg2008} --- для побудови та аналізу графів; matplotlib~\cite{Hunter2007} та seaborn --- для візуалізації; scikit-learn~\cite{Pedregosa2011} --- для обчислення NMI; python-louvain --- для алгоритму Лувена; scipy --- для обчислення кореляцій та KDE.

\subsection{Використання генеративного штучного інтелекту}
\label{sec:genai}

Відповідно до статті~8, частини~4 Закону України <<Про академічну доброчесність>>~\cite{AkademichnaDobrochesnist2024} (№~4742-IX від 2024~р., чинний з лютого 2025~р.), у роботі розкривається використання засобів генеративного штучного інтелекту.

У процесі підготовки курсової роботи використано Claude (Anthropic) --- велику мовну модель --- для таких цілей:

\begin{enumerate}
    \item \textbf{Допомога в написанні програмного коду:} генерація початкових версій Python-скриптів для бібліометричного аналізу, відлагодження коду, оптимізація алгоритмів. Усі скрипти перевірено та адаптовано автором.
    \item \textbf{Генерація чернеток тексту:} створення попередніх версій окремих розділів курсової роботи, які потім переглянуто, відредаговано та доповнено автором.
    \item \textbf{Аналіз бібліометричних даних:} інтерпретація кластерів VOSViewer, семантичний аналіз ключових слів, формулювання висновків за даними порівняльного аналізу.
    \item \textbf{Редагування та форматування:} LaTeX-верстка, виправлення стилістичних помилок, оформлення бібліографічних записів.
\end{enumerate}

Автор несе повну відповідальність за зміст роботи. Усі наукові висновки, інтерпретації результатів та методологічні рішення прийнято автором самостійно. Формальна декларація про використання ШІ наведена в Додатку~Г.
